
\section{Application Overview}

\section{Application Overview}

In this section, we describe the functionality of our application. The application can be started with several options:

\begin{itemize}
\item \textbf{Test Scene}: The scene to run, selected from a set of test scenes.
\item \textbf{Pipeline}: The rendering pipeline to use (see \autoref{subsec:pipelines}).
\item \textbf{Frame Count}: The number of frames to render.
\item \textbf{Resolution}: The resolution of the frames.
\item \textbf{SPP (Samples per Pixel)}: The number of rays cast per pixel (see \autoref{subsec:multi-sampling}).
\item \textbf{Delta Time}: The amount of animation time to advance between frames.
\item \textbf{Rebuild Frequency}: The number of frames after which the scene is rebuilt.
\item \textbf{Fast Build vs Fast Trace}: Whether to prioritize fast builds or fast ray tracing.
\end{itemize}

At startup, the application creates the specified ray tracing pipeline and loads the selected test scenes. It then applies the scene's animation for the current time, updates the nodes' global transforms, performs mesh skinning, and updates the scene's acceleration structures. After that, rays are traced, and the animation time is advanced by the delta time. These steps are repeated for the specified number of frames.

The application tracks the duration of acceleration structure updates, ray tracing, and their combined time for each frame. At the end of the run, it saves the tracked times along with their averages in a CSV file.

\subsection{Pipelines}\label{subsec:pipelines}

There are three possible pipelines:

\begin{itemize}
    \item \textbf{Basic Pipeline}
    
    This is the simplest pipeline, casting only primary rays. If a primary ray hits an object, the pixel is colored white; otherwise, it is colored black.

    \item \textbf{Standard Pipeline}
    
    This pipeline extends the basic pipeline. Its shaders have access to a directional light source and the scene's geometry, allowing for simple diffuse shading when a ray intersects an object.

    \item \textbf{Shadow Pipeline}
    
    This pipeline extends on the standard pipeline and uses secondary rays. When a primary ray hits a surface, a shadow ray is cast from the intersection point toward the light source. The ray checks if the path to the light is unobstructed before shading the sample.

\end{itemize}

\subsection{Multi-Sampling}\label{subsec:multi-sampling}

The application supports multi-sampling, which is a common anti-aliasing technique. We also use it to increase the number of rays in our test scenes without generating enormous images.

We use a straightforward method, where each pixel is divided into an $n \times n$ grid. This only works if the number of samples per pixel is a perfect square. One ray is cast through the center of each grid cell, and the final color of the pixel is determined by averaging all of the pixel's samples.

\begin{figure}[H]
    \centering
    \includegraphics[width=0.9\textwidth]{figures/multi_sampling.png}
    \caption{Multi-sampling grid}
    \label{fig:multi-sampling}
\end{figure}